\section[Spectral Algebras and Positivity]{Spectral Laplacian Algebras and a Uniform Positivity Criterion}
\label{sec:spectral-positivity}

From the spectral gradient-pairing table below and one Laplacian
coefficient bound, \cref{thm:spectral-gradient-subcalibration} constructs
\(G=F|\nabla F|\) for octonionic products and Pfaffians
(\cref{sec:norm-products-pfaffians}); complex determinants require
another multiplier (\cref{sec:complex-determinants}).

\subsection{Common gradient pairings and the smallest algebra}

Fix \(d\ge2\).  For \(r=(r_1,\ldots,r_d)\), let \(e_k(r)\)
be the elementary symmetric polynomial of degree \(k\), with
\(e_0=1\) and \(e_k=0\) for \(k<0\) or \(k>d\).
Consider real homogeneous polynomials \(F,E_1,\ldots,E_d\) on
\(\R^N\), of degrees \(d,2,\ldots,2d\), such that
\(F\not\equiv0\), \(E_1=R=|x|^2\), and at each point
\begin{equation}\label{eq:spectral-nonnegative-roots}
 (E_1(x),\ldots,E_d(x))=(e_1(r),\ldots,e_d(r))
 \quad\text{for some }r_1,\ldots,r_d\ge0.
\end{equation}
No differentiable choice or ordering of these roots is required.
Assume the identities
\begin{align}
 \Delta F&=0,& \Gamma(F,F)&=E_{d-1},
 \label{eq:spectral-F-pairings}\\
 \Gamma(F,E_k)&=2(d-k+1)E_{k-1}F,
 &&1\le k\le d,
 \label{eq:spectral-mixed-pairings}\\
 \Gamma(E_i,E_j)&=
 4\sum_{\alpha=1}^d r_\alpha
           \partial_\alpha e_i(r)\partial_\alpha e_j(r).
 \label{eq:spectral-root-pairings}
\end{align}
The right-hand side of \eqref{eq:spectral-root-pairings} is a
symmetric polynomial in the roots.  Its explicit polynomial form is
\begin{equation}\label{eq:spectral-polynomial-pairings}
 \Gamma(E_i,E_j)=4\sum_{\ell=0}^{i-1}
 (j-i+1+2\ell)E_{i-1-\ell}E_{j+\ell},
 \qquad 1\le i\le j\le d.
\end{equation}
To verify this formula, let \(p(t)=\prod_\alpha(1+r_\alpha t)\).
The generating function for one quarter of the pairings is
\[
 \sum_\alpha r_\alpha
       \partial_\alpha p(s)\partial_\alpha p(t)
 =\frac{st}{s-t}
       \bigl(sp'(s)p(t)-tp'(t)p(s)\bigr).
\]
The numerator vanishes at \(s=t\).  Polynomial division and
comparison of the coefficients of \(s^it^j\) give
\eqref{eq:spectral-polynomial-pairings}.

\begin{proposition}[Finite closure and exact generation]
\label{prop:spectral-smallest-algebra}
Suppose \eqref{eq:spectral-F-pairings}--\eqref{eq:spectral-root-pairings}
hold and
\begin{equation}\label{eq:spectral-Laplacian-drift}
 \Delta E_k=b_kE_{k-1}\qquad(1\le k\le d)
\end{equation}
for real constants \(b_k\).  Then
\(\R[E_1,\ldots,E_d,F]\) is a Laplacian algebra.
If \(d\ge3\) and \(b_k\ne0\) for \(2\le k\le d-1\), then
\begin{equation}\label{eq:spectral-smallest-algebra}
 \cA_R(F)=\R[E_1,\ldots,E_d,F].
\end{equation}
For \(d=2\), one instead has \(\cA_R(F)=\R[R,F]\).
If the \(d+1\) generators in \eqref{eq:spectral-smallest-algebra}
have independent differentials at some point, the spherical quotient
has dimension \(d\).
\end{proposition}

\begin{proof}
The generator criterion \cref{prop:MR-generator-test} gives closure.
Conversely, \(E_{d-1}=\Gamma(F,F)\in\cA_R(F)\), and successive
Laplacians recover \(E_{d-2},\ldots,E_1\) by
\eqref{eq:spectral-Laplacian-drift}.  For \(d\ge3\),
\eqref{eq:spectral-polynomial-pairings} gives
\begin{equation}\label{eq:norm-product-top-recovery}
 \Gamma(E_2,E_{d-1})
 =4\bigl((d-2)E_1E_{d-1}+dE_d\bigr),
\end{equation}
which recovers \(E_d\).  This proves equality of the algebras.
For \(d=2\), the identities
\(\Gamma(F,F)=R\), \(\Gamma(R,F)=4F\),
\(\Gamma(R,R)=4R\), \(\Delta F=0\), and \(\Delta R=2N\)
already close on \(R,F\).  Finally, independence gives quotient
dimension \(d+1\) in the Euclidean space; fixing the nonzero radial
generator \(E_1\) decreases this by one.
\end{proof}

The positivity theorem below requires only a lower bound for
\(\Delta E_{d-1}\), rather than the full drift formula
\eqref{eq:spectral-Laplacian-drift}.

\subsection{Two elementary inequalities}

Set
\begin{equation}\label{eq:spectral-SBCD}
 S=E_{d-1},\qquad B=E_{d-2},\qquad
 C=E_{d-3},\qquad D=E_d.
\end{equation}
In this subsection the letter \(C\) denotes this invariant, not the
zero cone.  All formulas will be used on \(S>0\), where
\(B>0\) and every \(E_k\) with \(0\le k\le d-1\) is positive.

\begin{lemma}[A Newton inequality and an antitone coefficient]
\label{lem:spectral-Newton-monotonicity}
For nonnegative roots and \(d\ge3\),
\begin{equation}\label{eq:spectral-Newton}
 BS\ge\frac{3d}{d-2}CD.
\end{equation}
For every \(d\ge2\), the function
\begin{equation}\label{eq:spectral-Q}
 Q_* =\frac{5B}{S}+\frac{3CD}{S^2}
\end{equation}
is positive and nonincreasing in each root \(r_i\) wherever
\(S>0\).
\end{lemma}

\begin{proof}
First suppose all roots are positive and put \(t_i=r_i^{-1}\).
Writing \(p_k=e_k(t_1,\ldots,t_d)\), one has
\(S=Dp_1\), \(B=Dp_2\), and \(C=Dp_3\).
For each triple of nonnegative numbers, the sum of the six monomials
\(t_i^2t_j\), with distinct indices in the triple, is at least
\(6t_it_jt_k\), by the arithmetic--geometric mean inequality.
Summing over triples counts each ordered pair \(d-2\) times.
Consequently
\[
 p_1p_2=\sum_{i\ne j}t_i^2t_j+3p_3
 \ge\frac{3d}{d-2}p_3,
\]
which proves \eqref{eq:spectral-Newton}; continuity includes zero roots.

In the reciprocal variables, \(Q_* =5p_2/p_1+3p_3/p_1^2\).
Fix \(t_i=z\), and let \(u,v,w\) be the elementary symmetric
polynomials of orders \(1,2,3\) in the remaining variables.
Then \(p_1=z+u\), \(p_2=zu+v\), and \(p_3=zv+w\), so
\begin{equation}\label{eq:spectral-Q-derivative}
 \frac{\partial Q_*}{\partial z}
 =\frac{z(5u^2-8v)+5u^3-2uv-6w}{(z+u)^3}\ge0.
\end{equation}
Indeed \(u^2\ge2v\) and \(u^3\ge6w\); the numerator is at least
\(zu^2+3u^3\).  Since \(\partial_{r_i}=-t_i^2\partial_{t_i}\),
this proves the asserted sign.  Formula \eqref{eq:spectral-Q}
is smooth in the original roots wherever \(S>0\), so its derivative
sign extends to zero roots by continuity.  Positivity follows from
\(B>0\).  For \(d=2\), the same argument has \(v=w=0\), or
one can use \(Q_*=5/E_1\) directly.
\end{proof}

\subsection{Subcalibration existence from a Laplacian bound}

\begin{theorem}[Spectral subcalibrations]
\label{thm:spectral-gradient-subcalibration}
Assume \eqref{eq:spectral-nonnegative-roots}--\eqref{eq:spectral-root-pairings}.
Suppose on \(U=\{S>0\}\) that
\begin{equation}\label{eq:spectral-threshold}
 \Delta S\ge\beta B,\qquad
 \beta\ge24-\frac{16}{d}.
\end{equation}
Then \(N\ge4d\), and there is a potential \(G\), homogeneous of
degree \(2d-1\),
with nonzero gradient on \(U\) and satisfying \(F\cL(G)\ge0\).
It is given explicitly by
\begin{equation}\label{eq:spectral-gradient-potential}
 G=F\sqrt S=F|\nabla F|.
\end{equation}
Writing \(X=\nabla G/|\nabla G|\), one moreover has
\begin{equation}\label{eq:spectral-strict-divergence}
 \operatorname{sgn}(F)\diver X>0
 \qquad\text{on }\{F\ne0\}.
\end{equation}
The field \(X\) is a two-sided subcalibration of \(\{F<0\}\)
across \(Z=\{S=0\}\), and the perimeter and current conclusions of
\cref{thm:removable-subcalibration} hold. The cone
\(\Sigma=\{F=0\}\) is strictly area minimizing and strictly
stable, with
\begin{equation}\label{eq:spectral-Hardy-gap}
 Q_\Sigma(u)\ge\frac{(N-4d+1)^2}{4}
       \int_\Sigma r^{-2}u^2,
 \qquad u\in C_c^\infty(\Sigma_{\reg}\setminus\{0\}).
\end{equation}
The constant in this estimate is at least \(1/4\).
\end{theorem}

\begin{proof}
For positive roots, the decreasing ratios of the normalized elementary symmetric
functions \(e_j/\binom dj\), as given by Newton's inequalities
\cite[Section~3.1]{BorceaBrandenLiggett09}, imply
\begin{equation}\label{eq:spectral-radial-Newton}
 R B\ge\frac{d(d-1)}2S.
\end{equation}
This also holds at zero roots by continuity, and is an equality for
\(d=2\). Since \(S=|\nabla F|^2\) is a nonzero polynomial,
\(U\) is dense, so \eqref{eq:spectral-threshold} extends to all
points by continuity. Integrating on the unit sphere and using the
homogeneous degree \(2d-2\) of \(S\) gives
\[
 (2d-2)(N+2d-4)\int_{\Sph^{N-1}}S
 =\int_{\Sph^{N-1}}\Delta S
 \ge(12d-8)(d-1)\int_{\Sph^{N-1}}S.
\]
The integral of \(S\) is positive. Division proves \(N\ge4d\).

The mixed pairing at \(k=d-1\) gives
\(\Gamma(F,S)=4BF\).  Thus the regular zero set is minimal and
\eqref{eq:harmonic-minimal-weight} holds with \(q=2B\).
By \eqref{eq:spectral-polynomial-pairings}, \(T=\Gamma(S,S)=4(BS+3CD)\).
Take \(h=\tfrac12\log S\), and write
\(\delta=\beta/4\).  The auxiliary function in
\eqref{eq:harmonic-gradient-multiplier-data} is exactly \(Q_*\).
Moreover,
\begin{align}
 \Delta h&\ge2(\delta-1)\frac BS-6\frac{CD}{S^2},
 \label{eq:spectral-log-Laplacian}\\
 D_h&\ge2(\delta-4)B-12\frac{CD}{S}.
 \label{eq:spectral-boundary-bound}
\end{align}
For \(d\ge3\), \cref{lem:spectral-Newton-monotonicity} implies
\(BS\ge3CD\); hence
\(\Delta h\ge2(\delta-2)B/S>0\).  The sharper constant gives
\[
 D_h\ge2\left(\delta-6+\frac4d\right)B\ge0.
\]
For \(d=2\), \(C=0\) and \(\delta\ge4\), so
\(\Delta h\ge2(\delta-1)/S>0\) and \(D_h\ge0\) follow
directly from \eqref{eq:spectral-log-Laplacian}--\eqref{eq:spectral-boundary-bound}.

For every smooth symmetric function \(u(r)\)
expressed in the generators, the mixed pairings give
\begin{equation}\label{eq:spectral-root-derivation}
 \Gamma(F,u)=2F\sum_i\partial_{r_i}u,
\end{equation}
because \(\sum_i\partial_{r_i}e_k=(d-k+1)e_{k-1}\).
The root pairing similarly gives
\[
 \Gamma(h,Q_*)=
 \frac2S\sum_i r_i(\partial_{r_i}S)(\partial_{r_i}Q_*)\le0,
 \qquad
 F\Gamma(F,Q_*)=2F^2\sum_i\partial_{r_i}Q_*\le0.
\]
Here \(\partial_{r_i}S\ge0\), and the other derivative has the
opposite sign by \cref{lem:spectral-Newton-monotonicity}.
The invariant identities extend to repeated and zero roots, verifying
\cref{thm:harmonic-monotonicity-criterion}. In particular,
\begin{equation}\label{eq:spectral-gradient-norm}
 |\nabla G|^2=S^2+5BF^2+\frac{3CD}{S}F^2\ge S^2>0.
\end{equation}
Since \(Q_*>0\) and \(\Delta h>0\), the term
\(F^2Q_*\Delta h\) in \eqref{eq:harmonic-log-defect} is strictly
positive when \(F\ne0\), while all the other terms in that bracket
are nonnegative. Thus \(F\cL(G)>0\) there. Euler's identity shows
that \(F\ne0\) implies \(S>0\), so this proves
\eqref{eq:spectral-strict-divergence} on both entire phases.

Harmonicity makes \(Z\) removable by
\cref{lem:square-free-critical-set}; the monotonicity theorem gives
comparison with full support by \cref{rem:polynomial-finite-perimeter}.
Both phases are nonempty, and \eqref{eq:spectral-strict-divergence}
verifies the flux hypothesis of \cref{thm:strict-minimality-flux}.
Homogeneity therefore gives strict area minimality.

For strict stability, we use the boundary Jacobi estimate.
At a smooth point of the full harmonic
nodal hypersurface, the
two local phases have opposite signs, and the Hopf boundary lemma
gives \(\nabla F\ne0\). Thus the regular cone lies in \(U\), where
\eqref{eq:harmonic-log-Jacobi} gives
\(V_G=D_h/S\ge0\). Apply \cref{thm:ground-state-gap} with
\(n=N-1\) and \(\ell=2d-1\) to obtain
\eqref{eq:spectral-Hardy-gap}, since
\[
 \left(\ell-\frac n2\right)^2
 =\frac{(N-4d+1)^2}{4}.
\]
The bound \(N\ge4d\) makes this constant at least \(1/4\).
\end{proof}

\begin{remark}[An explicit coefficient identity]
\label{rem:spectral-explicit-certificate}
If \(\Delta S=4\delta B\) exactly, put
\(A=E_{d-4}\) and \(J=E_{d-5}\), retaining the convention for
negative indices.  The same closure table gives the checkable identity
\begin{equation}\label{eq:spectral-explicit-certificate}
 \cL(F\sqrt S)=\frac{F}{2S^{3/2}}
 \left\{4S^2\bigl[(\delta-4)BS-6CD\bigr]+F^2\mathscr P_\delta\right\},
\end{equation}
where
\[
 \mathscr P_\delta=(20\delta+10)B^2S+(12\delta-18)BCD
                    -56CS^2-82ASD-30JD^2.
\]
\Cref{thm:spectral-gradient-subcalibration} proves the sign of this
identity by monotonicity, uniformly in \(d\).
\end{remark}

\subsection{The threshold and the role of the Laplacian}

The threshold is sharp for \(G=F\sqrt S\) when
\(\Delta S=4\delta B\) and the invariant image contains
\(r_1=\cdots=r_d=1\) with arbitrarily small nonzero \(F\).
Indeed, at these roots,
\begin{equation}\label{eq:spectral-balanced-defect}
 D_h=d(d-1)\left(\delta-6+\frac4d\right).
\end{equation}
For \(\delta<6-4/d\), the multiplier sign fails near \(F=0\).
Complex determinants realize these values with \(\delta=4\), so
\(F\sqrt S\) fails in every order \(d\ge3\). Their product multiplier
in \cref{sec:complex-determinants} succeeds: the threshold is not
necessary for area minimization.

The next proposition compares reduced operators with different
Laplacian coefficients; their ambient spaces may differ.

\begin{proposition}[Monotonicity under an increase of the drift]
\label{prop:spectral-drift-comparison}
Suppose two reduced operators have the same gradient pairing
\(\widehat\Gamma\), annihilate \(F\), and have generator drifts
\(\ell_i=\widehat\Delta E_i\) and
\(\widetilde\ell_i=\widetilde\Delta E_i\).  These are the full
drift functions, including any factors \(E_{i-1}\).
For a positive smooth multiplier
\(a=a(E_1,\ldots,E_d)\), put \(G=aF\) and
\(W=\widehat\Gamma(G,G)\).  Their curvature operators satisfy
\begin{equation}\label{eq:spectral-drift-comparison}
 G\widetilde{\cL}(G)-G\cL(G)
 =F^2aW\sum_i(\widetilde\ell_i-\ell_i)\partial_{E_i}a.
\end{equation}
On any common invariant domain where the reference certificate holds,
a nonnegative sign certificate is preserved if
\(\widetilde\ell_i\ge\ell_i\) for every generator used by \(a\)
and \(\partial_{E_i}a\ge0\).
\end{proposition}

\begin{proof}
The common gradient pairing gives the same \(W\) and the same
term \(-\widehat\Gamma(G,W)/2\).  The difference of the two
Laplacians on \(G\) is
\(F\sum_i(\widetilde\ell_i-\ell_i)\partial_{E_i}a\).
Multiplication by \(GW\) proves the formula.
\end{proof}

For example, consider drifts of the form
\begin{equation}\label{eq:spectral-two-parameter-drift}
 \Delta E_k=2j[\alpha+\eta(j-1)]E_{k-1},
 \qquad j=d-k+1.
\end{equation}
The simple criterion requires \(\alpha+\eta\ge6-4/d\).
The determinant drift \((\alpha,\eta)=(2,2)\) instead admits
\(a=(E_1\cdots E_{d-1})^{1/2}\) by
\cref{sec:complex-determinants}. Its nonnegative partial derivatives
allow \cref{prop:spectral-drift-comparison} to transfer the certificate
whenever
\begin{equation}\label{eq:spectral-determinant-drift-range}
 \alpha+\eta\ge4,\qquad
 \alpha+(d-1)\eta\ge2d.
\end{equation}
Indeed these endpoint conditions are equivalent to
\(\alpha+\eta(j-1)\ge2j\) for \(2\le j\le d\).
To apply this comparison geometrically, the field must be defined off
a removable exceptional set.
